\songtitle{Un día}

\tag*{2026}\tag{trova}\tag{spanish-lyrics}
\begin{notes}
Lyrics by SUNO based on a short story by Claude written as an exercise of making sense of a series of photographs by Carlos Thompson as a small urban adventure.
\end{notes}
\begin{style}
Nueva trova cubana with gentle 6/8 guitar pulse, light percussion and bass, verse-led storytelling that breathes between images; pre-chorus lifts with held harmony and a tumbao hint, chorus opens with warm unison vocals and short call-and-response, bridge turns reflective and sparse before a final chorus with fuller backing voices, Intimate close-mic lead, soft doubles on key lines, subtle room echo, brief acoustic fills and bird-like guitar harmonics between phrases, warm organic mix with earthy mids, trova, simple
\end{style}
\begin{lyrics}
\songpart{Verse 1}
Cada mañana igual,\\
Mochi en su planeta.\\
Chocolate y calor\\
sobre la cama blanca.

Yo lo miro en la puerta,\\
sonrío sin hablar.\\
El espresso ya espera,\\
negro, fino, total.

\songpart{Pre-Chorus}
Combustible del día,\\
mi pequeño ritual.\\
Y afuera la ciudad\\
ya empieza a empujar.

\songpart{Chorus}
Un día, un día,\\
todo vuelve a empezar.\\
Un día, un día,\\
la llave y el café.\\
Un día, un día,\\
Bogotá me deja ver\\
que entre concreto y niebla\\
todavía nace la fe.

\songpart{Verse 2}
Bogotá no perdona,\\
se desborda el motor.\\
Por la ventanilla\\
van motos sin control.

Placas amarillas,\\
río de metal.\\
Corre todo el mundo,\\
nadie llega al final.

\songpart{Pre-Chorus}
Pero arriba, entre el humo,\\
las montañas están.\\
Y guardan algo vivo\\
detrás del capital.

\songpart{Chorus}
Un día, un día,\\
todo vuelve a empezar.\\
Un día, un día,\\
la llave y el café.\\
Un día, un día,\\
Bogotá me deja ver\\
que entre concreto y niebla\\
todavía nace la fe.

\songpart{Verse 3}
Mi HP me espera\\
sobre el frío del metal.\\
Mesa de coworking,\\
tejados a mirar.

El profe en la pantalla:\\
fusibles, regulador,\\
condensadores, futuro,\\
trabajo y motor.

Fotografío el mapa,\\
me lo quiero aprender.\\
La ciudad me habla\\
sin necesidad de ver.

\songpart{Verse 4}
En la pausa cae rojo\\
un tarro junto al cristal.\\
Afuera, los árboles\\
brillan como un vendaval.

La llave turquesa\\
sobre la mesa oscura\\
me dice que regreso,\\
que la casa me procura.

\songpart{Bridge}
Y en la pantalla púrpura\\
alguien me hizo reír:\\
si no amas la fresa,\\
toma chocolate y sé feliz.

Porque así de simple\\
debería ser vivir:\\
entender al otro,\\
dejarlo compartir.

\songpart{Verse 5}
Vi una mujer indígena\\
con dos golondrinas más.\\
Tenía los ojos cerrados,\\
como escuchando la ciudad.

Me quedé largo rato,\\
sin moverme, sin pensar.\\
Bogotá también sueña,\\
la vi respirar.

\songpart{Chorus}
Un día, un día,\\
todo vuelve a empezar.\\
Un día, un día,\\
la llave y el café.\\
Un día, un día,\\
Bogotá me deja ver\\
que entre concreto y niebla\\
todavía nace la fe.

\songpart{Outro}
Al final abro la puerta,\\
Mochi vuelve a esperar.\\
Mismo punto exacto,\\
mismo círculo lunar.

Me siento, sirvo otro,\\
y comienzo a escribir:\\
quizás no pasó el tiempo,\\
quizás aprendí a vivir.
\end{lyrics}
